\documentclass[11pt, aspectratio=169]{beamer}

\usetheme{default}
\usefonttheme{serif}
\setbeamertemplate{navigation symbols}{}
\setbeamercolor{normal text}{fg=black, bg=white}
\setbeamercolor{structure}{fg=black}
\setbeamercolor{frametitle}{fg=black, bg=white}
\setbeamercolor{title}{fg=black}
\setbeamertemplate{frametitle}{%
  \vspace{3mm}{\large\bfseries\insertframetitle}\par\vspace{1.2mm}%
  \hrule height 0.5pt\vspace{2.5mm}}
\setbeamertemplate{footline}{%
  \hfill{\footnotesize\insertframenumber\,/\,\inserttotalframenumber}\hspace{4mm}\vspace{2.5mm}}
\setbeamertemplate{itemize item}{\rule[0.5ex]{1.1ex}{0.5pt}}
\setbeamertemplate{itemize subitem}{\rule[0.5ex]{0.8ex}{0.5pt}}
\setbeamersize{text margin left=9mm, text margin right=9mm}

\usepackage{amsmath, amssymb}
\usepackage{booktabs}
\usepackage{listings}
\usepackage[most]{tcolorbox}
\usepackage{circuitikz}
\usetikzlibrary{patterns}

\lstdefinelanguage{Rust}{
  keywords={fn, let, mut, if, else, while, for, in, match, return, struct,
    enum, impl, trait, pub, use, const, loop, continue, break, ref, type,
    where, as, self, Self, true, false, Some, None, Ok, Err},
  sensitive=true,
  comment=[l]{//},
  string=[b]",
}

\newtcblisting{rust}{listing only, sharp corners, colback=white,
  colframe=black, boxrule=0.5pt, left=5pt, right=4pt, top=2pt, bottom=2pt,
  listing options={language=Rust, basicstyle=\footnotesize\ttfamily,
  commentstyle=\itshape, keywordstyle=\bfseries,
  columns=fullflexible, keepspaces=true, breaklines=true}}

\newtcolorbox{tile}[1]{sharp corners, colback=white, colframe=black,
  colbacktitle=white, coltitle=black, titlerule=0.5pt,
  fonttitle=\bfseries\small, title={#1}, boxrule=0.5pt,
  left=6pt, right=6pt, top=4pt, bottom=5pt, fontupper=\small}

% loc tile: code section name, label, budget range [start, end] of 992
\newcommand{\locbar}[4]{%
\begin{tcolorbox}[sharp corners, colback=white, colframe=black,
  boxrule=0.5pt, left=6pt, right=6pt, top=3pt, bottom=5pt]
{\footnotesize\texttt{#1}\hfill #2}\\[3pt]
\begin{tikzpicture}
\pgfmathsetlengthmacro{\bw}{\linewidth-4pt}
\fill[pattern=north east lines] (0,0) rectangle ({(#3/992)*\bw}, 4pt);
\fill[black] ({(#3/992)*\bw},0) rectangle ({(#4/992)*\bw}, 4pt);
\draw[black, line width=0.4pt] (0,0) rectangle (\bw, 4pt);
\end{tikzpicture}%
\end{tcolorbox}}

\title{\textbf{nanospice}}
\subtitle{A SPICE circuit simulator in 1000 lines of Rust}
\author{Milan Rother}
\date{2026 \\[3mm] {\footnotesize github.com/milanofthe/nanospice}}

\begin{document}

\begin{frame}
\titlepage
\end{frame}

\begin{frame}{The rule}
\begin{itemize}
\item one file, \texttt{src/main.rs}, at most 1000 nonblank, noncomment lines
\item standard library only, no dependencies
\item a test counts the lines and fails the build above the limit
\item tests, comments, and example netlists are free
\end{itemize}
\vspace{3mm}
The budget forces every feature to displace another, which makes the cost
of each decision explicit, and it keeps the simulator readable in one
sitting.
\vspace{3mm}
\begin{center}
current count: \textbf{992 of 1000}
\end{center}
\end{frame}

\begin{frame}{What fits}
\begin{columns}[T]
\column{0.48\textwidth}
\textbf{Analyses}
\begin{itemize}
\item \texttt{.op}, \texttt{.dc} sweep
\item \texttt{.tran} with adaptive timestep, \texttt{uic}
\item \texttt{.ac} dec/lin
\end{itemize}
\vspace{2mm}
\textbf{Devices}
\begin{itemize}
\item R, C, L, V, I, E, G
\item diode, MOSFET, JFET, BJT
\item junction and gate capacitances
\end{itemize}
\column{0.48\textwidth}
\textbf{Solver}
\begin{itemize}
\item sparse LU with partial pivoting
\item Newton with junction limiting
\item gmin and source stepping
\item trapezoid, LTE control, breakpoints
\end{itemize}
\vspace{2mm}
\textbf{Input}
\begin{itemize}
\item SIN, PULSE, PWL waveforms
\item \texttt{.model}, \texttt{.print} cards
\end{itemize}
\end{columns}
\vspace{3mm}
\begin{center}
23 integration tests against analytic references, none of them counted
\end{center}
\end{frame}

\begin{frame}{Where the lines go}
\begin{center}
\begin{tabular}{lrr}
\toprule
code section & lines & cumulative \\
\midrule
constants, error and output plumbing & 23 & 23 \\
numbers: \texttt{Num} trait, complex, unit suffixes & 78 & 101 \\
circuit data model & 61 & 162 \\
parser & 275 & 437 \\
sparse LU and stamp helpers & 113 & 550 \\
device models & 64 & 614 \\
solver: assembly, Newton, operating point & 142 & 756 \\
output selection & 26 & 782 \\
analyses and main & 210 & 992 \\
\bottomrule
\end{tabular}
\end{center}
\vspace{2mm}
The parser is the largest section: the numerics of a circuit simulator
are compact, the input format is not.
\end{frame}

\begin{frame}{Modified nodal analysis}
Unknowns: node voltages plus one branch current per voltage-defined
element (V, L, E). KCL per node, one constraint row per branch.
\vspace{2mm}
\begin{columns}[c]
\column{0.4\textwidth}
\begin{center}
\begin{circuitikz}[american, scale=0.75, transform shape]
\draw (0,0) node[ground]{} to[V, invert, l=$10\,\mathrm{V}$] (0,2.4) -- (1.8,2.4)
  node[circ]{} node[above]{$v_1$}
  to[R, l=$1\,\mathrm{k\Omega}$] (4.4,2.4) node[circ]{} node[above]{$v_2$}
  to[R, l=$1\,\mathrm{k\Omega}$] (4.4,0) node[ground]{};
\end{circuitikz}
\end{center}
\column{0.58\textwidth}
\[
\begin{pmatrix}
g & -g & 1 \\
-g & 2g & 0 \\
1 & 0 & 0
\end{pmatrix}
\begin{pmatrix} v_1 \\ v_2 \\ i \end{pmatrix}
=
\begin{pmatrix} 0 \\ 0 \\ 10 \end{pmatrix}
\]
\[
v_1 = 10, \quad v_2 = 5, \quad i = -5\,\text{mA}
\]
\end{columns}
\vspace{2mm}
\begin{itemize}
\item the branch row has a structural zero on the diagonal: pivoting is mandatory
\item $i < 0$: branch current flows into the positive terminal, a sourcing supply reads negative
\end{itemize}
\end{frame}

\begin{frame}[fragile]{The ground row}
\begin{tile}{Design decision: a dummy ground row}
Ground is unknown 0 and the solver loops start at 1. The alternative is
an \emph{if node is not ground} test inside every stamp, roughly two
dozen branches across eleven devices. One guard inside \texttt{add}
replaces all of them; $x_0 = 0$ is pinned after the solve.
\end{tile}
\vspace{2mm}
\begin{rust}
fn quad(&mut self, p: usize, n: usize, cp: usize, cn: usize, g: T) {
    self.add(p, cp, g);
    self.add(p, cn, T::zero() - g);
    self.add(n, cp, T::zero() - g);
    self.add(n, cn, g);
}
\end{rust}
The conductance quad covers R and the VCCS; \texttt{branch} covers V, L,
and E. Every stamp in the file reduces to these two shapes.
\end{frame}

\begin{frame}{Sparse LU: the data structure}
\locbar{sparse LU and stamp helpers}{113 lines, budget 437 to 550}{437}{550}
\vspace{2mm}
\begin{itemize}
\item each matrix row: a vector of (column, value) pairs, sorted by column
\item eliminated columns are removed from active rows
\end{itemize}
\vspace{2mm}
\begin{tile}{The first-entry invariant}
Before elimination step $k$, every row stored at index $i \ge k$ contains
only columns $\ge k$. Consequence: if such a row has an entry in column
$k$ at all, it is the \emph{first} element of its vector. Pivot search is
an $O(1)$ probe per candidate row, with no column-linked index structure.
\end{tile}
\end{frame}

\begin{frame}[fragile]{Sparse LU: elimination is a merge}
\begin{rust}
// row a minus f times pivot row p, both sorted
fn merge<T: Num>(a: &[(usize, T)], p: &[(usize, T)], f: T) -> Vec<(usize, T)> {
    let (mut i, mut j) = (0, 0);
    let mut o = Vec::with_capacity(a.len() + p.len());
    while i < a.len() && j < p.len() { /* three-way column compare */ }
    // fill-in falls out of the merge as the union of the column sets
    o
}
\end{rust}
\vspace{1mm}
\begin{itemize}
\item pivoting: largest magnitude among the candidate first entries, rows swapped whole
\item no candidate for column $k$: matrix is singular, the failing column is named \\
      {\footnotesize\texttt{singular matrix at node b, no conduction path?}}
\item no Markowitz ordering: netlist order, linear cost on ladders and trees
\item 2000-node RC ladder, OP plus AC sweep: tens of milliseconds
\end{itemize}
\end{frame}

\begin{frame}[fragile]{One primitive for every nonlinear device}
A static device model is a set of current branches, each controlled by
node-pair voltages. This is the contribution operator of Verilog-A,
$I(p,n) \mathrel{{<}{+}} f(v)$:
\begin{rust}
fn contrib(&mut self, p: usize, n: usize, i: T,
           ctl: &[(usize, usize, T)], x: &[T]) {
    let mut ieq = i;
    for &(a, b, g) in ctl {
        self.quad(p, n, a, b, g);
        ieq = ieq - g * (x[a] - x[b]);
    }
    self.rhs(p, T::zero() - ieq);
    self.rhs(n, ieq);
}
\end{rust}
The dual \texttt{vcontrib} writes the branch row
$v_p - v_n - \sum_k g_k (x_{a_k} - x_{b_k}) = v$ for the
voltage-defined elements.
\end{frame}

\begin{frame}{The whole device set on two primitives}
\begin{center}
\begin{tabular}{lll}
\toprule
device & primitive & controls \\
\midrule
R, C, I & \texttt{contrib} & self conductance, none for I \\
G (VCCS) & \texttt{contrib} & one control pair \\
diode & \texttt{contrib} & one junction \\
MOSFET, JFET & \texttt{contrib} & $g_m$ and $g_{ds}$ \\
BJT & 3 $\times$ \texttt{contrib} & two junctions, transport source \\
V & \texttt{vcontrib} & none \\
L & \texttt{vcontrib} & its own branch current \\
E (VCVS) & \texttt{vcontrib} & one control pair \\
\bottomrule
\end{tabular}
\end{center}
\vspace{2mm}
The refactor to this form removed more lines than the primitives cost;
the freed budget paid for the JFET and the PWL source. Hand-derived
Jacobian stamps become verification artifacts.
\end{frame}

\begin{frame}{Diode and junction limiting}
Newton diverges on the bare exponential: an early iterate proposing
$5\,$V puts $e^{193}$ into the Jacobian.
\[
v_{crit} = V_T \ln \frac{V_T}{\sqrt{2}\, I_S}
\quad \text{(the point of maximum curvature of } i(v)\text{)}
\]
Beyond $v_{crit}$, \texttt{pnjlim} replaces the proposed step by a
logarithmic one:
\[
v \leftarrow v_{old} + V_T \ln\!\left(1 + \frac{v_{new} - v_{old}}{V_T}\right)
\]
\vspace{1mm}
\begin{itemize}
\item the junction advances a controlled amount per iteration, regardless of overshoot
\item per-device memory of the last limited voltage; the BJT carries two
\end{itemize}
\end{frame}

\begin{frame}{Limiting must deny convergence}
\begin{tile}{Field note}
With limiting alone, a common-emitter amplifier converged with the
transistor off: the divider put the base at 2.1\,V, \texttt{pnjlim}
clamped the junction seen by the stamps to 0.11\,V, no current flowed,
the node voltages stopped changing, and the delta test declared success.
The output was superficially plausible: the divider voltage was correct,
but unloaded.
\end{tile}
\vspace{2mm}
\begin{itemize}
\item repair: an iteration in which a limiter changed a junction voltage cannot converge
\item the assembly reports a \texttt{limited} flag, the Newton loop requires it false
\item with the flag: correct bias (2.05\,V, 1.3\,mA) in a dozen iterations
\item a convergence test on the solution vector alone cannot detect active limiting
\end{itemize}
\end{frame}

\begin{frame}{MOSFET: the effective frame}
Four cases (nmos/pmos, source-drain swap), one evaluation. With polarity
$s = \pm 1$ and effective terminals chosen so $v_{DS}^e \ge 0$:
\[
\frac{\partial I}{\partial v_g}
 = s \cdot g_m \cdot s = g_m
\]
Every derivative picks up the polarity twice, and $s^2 = 1$: the
conductance stamps are identical for all four cases, only the right-hand
side carries $s$.
\vspace{2mm}
\begin{itemize}
\item square law with channel-length modulation, $g_m$ and $g_{ds}$ by differentiation
\item the JFET is the same square law with $K_P = 2\beta$ and a depletion
      threshold: it shares the MOSFET variant and parser arm
\end{itemize}
\end{frame}

\begin{frame}[fragile]{BJT: three contributions}
\begin{columns}[c]
\column{0.4\textwidth}
\begin{center}
\begin{circuitikz}[american, scale=0.72, transform shape]
\draw (2.6,4.4) node[circ]{} node[above]{$c$} -- (2.6,3.6);
\draw (0,2.2) node[circ]{} node[left]{$b$} -- (0.7,2.2);
\draw (0.7,2.2) to[D, l=$i_{BC}$] (2.6,3.6);
\draw (0.7,2.2) to[D, l_=$i_{BE}$] (2.6,0.8);
\draw (2.6,0.8) -- (2.6,0) node[circ]{} node[below]{$e$};
\draw (2.6,4.4) -- (4.8,4.4) to[cI, l=$i_{CT}$] (4.8,0) -- (2.6,0);
\end{circuitikz}
\end{center}
\column{0.58\textwidth}
\[
i_{CT} = I_S (e_f - e_r), \;\;
i_{BE} = \tfrac{I_S}{\beta_F}(e_f - 1), \;\;
i_{BC} = \tfrac{I_S}{\beta_R}(e_r - 1)
\]
\begin{rust}
s.contrib(*b, *e, .., &[(*b, *e, gpi)], x);
s.contrib(*b, *c, .., &[(*b, *c, gmu)], x);
s.contrib(*c, *e, ..,
    &[(*b, *e, gmf), (*b, *c, -gmr)], x);
\end{rust}
\end{columns}
\vspace{2mm}
\begin{itemize}
\item KCL composes the classic $3 \times 3$ stamp; rows and columns sum to
      zero by construction
\item pnp through the same sign argument; an npn/pnp mirror test requires
      exact symmetry
\end{itemize}
\end{frame}

\begin{frame}{Junction capacitances by desugaring}
A capacitance across a node pair is exactly the \texttt{C} device, so the
parser desugars \texttt{cjo}, \texttt{cje}/\texttt{cjc},
\texttt{cgs}/\texttt{cgd} into internal capacitors. No assembly, state,
timestep, or AC code knows they exist.
\[
C(v) = C_{j0} \left(1 - \tfrac{v}{V_J}\right)^{-m}, \qquad
Q(v) = \frac{C_{j0} V_J}{1 - m}
       \left(1 - \left(1 - \tfrac{v}{V_J}\right)^{1 - m}\right)
\]
\begin{itemize}
\item charge-based capacitor, linearized above $F_C V_J$; $m = 0$ is the plain linear C
\item junctions graded ($m = 0.5$), gate oxide linear, \texttt{m=} available on the C card
\item test: at $v = -3$\,V the capacitance is exactly $C_{j0}/2$, the AC corner doubles
\end{itemize}
\end{frame}

\begin{frame}[fragile]{Operating point: homotopy paths as data}
\locbar{solver: assembly, Newton, operating point}{142 lines, budget 614 to 756}{614}{756}
\vspace{1mm}
Both classic fallbacks are paths in the $(g_{min}, \text{source scale})$
plane, walked with warm starts:
\begin{rust}
let gmin_path: Vec<(f64, f64)> =
    (2..=12).map(|k| (10f64.powi(-k), 1.0)).collect();
let src_path: Vec<(f64, f64)> =
    (1..=10).map(|k| (GMIN, k as f64 / 10.0)).collect();
for path in [gmin_path, src_path] {
    if path.iter().all(|&(g, fac)|
        newton(ckt, x, lim, &Mode::Dc { fac }, g, ITL_OP).is_some()) {
        return;
    }
}
\end{rust}
\begin{itemize}
\item gmin stepping: junction trouble; source stepping: feedback and bistable topologies
\item a cross-coupled NMOS latch in the test suite exercises the path
\end{itemize}
\end{frame}

\begin{frame}{Transient: companion models}
Trapezoidal rule on $i = C\,dv/dt$: the capacitor becomes a resistor with
a history source for one step. The inductor is the dual through the flux
$\varphi = L i$.
\begin{center}
\begin{circuitikz}[american, scale=0.62, transform shape]
\draw (0,3.4) node[circ]{} node[above]{$p$} to[C, l=$C$] (0,0) node[circ]{} node[below]{$n$};
\node at (1.5,1.7) {$\rightarrow$};
\draw (3.0,3.4) node[circ]{} node[above]{$p$} -- (5.4,3.4);
\draw (3.0,3.4) to[R, l=$\tfrac{2C}{h}$] (3.0,0);
\draw (5.4,3.4) to[I, l=$I_{eq}$] (5.4,0);
\draw (3.0,0) node[circ]{} node[below]{$n$} -- (5.4,0);
\end{circuitikz}
\hspace{10mm}
\begin{circuitikz}[american, scale=0.62, transform shape]
\draw (0,3.4) node[circ]{} node[above]{$p$} to[L, l=$L$] (0,0) node[circ]{} node[below]{$n$};
\node at (1.5,1.7) {$\rightarrow$};
\draw (3.0,3.4) node[circ]{} node[above]{$p$} to[R, l=$\tfrac{2L}{h}$] (3.0,1.7) to[V, l=$V_{eq}$] (3.0,0) node[circ]{} node[below]{$n$};
\end{circuitikz}
\end{center}
\begin{itemize}
\item state per reactive device: charge or flux plus the dual rate, one update rule for both
\item trapezoid: second order, A-stable, no numerical damping (an LC tank
      holds its amplitude for twenty periods, a test enforces it)
\end{itemize}
\end{frame}

\begin{frame}{Timestep control by truncation error}
\[
\varepsilon_{tr} = -\tfrac{h^3}{12} \dddot{x}, \qquad
\varepsilon_{tr} \approx
\frac{h^3/12}{\;h(h+d_1)(h+d_1+d_2)/6 + h^3/12\;}\,\bigl(x_c - x_p\bigr)
\]
\begin{itemize}
\item quadratic predictor through the last three accepted points, Lagrange
      coefficients for variable steps
\item corrector minus predictor measures $\dddot{x}$ (Milne estimate);
      uniform steps give the classic factor $1/13$
\item cube-root step rule $h' = 0.9\, h\, r^{-1/3}$, clamped to $[0.3, 2] \times h$
\item the predictor doubles as the Newton starting point
\item nonconvergence: cut the step by 8, the SPICE2 default
\end{itemize}
\end{frame}

\begin{frame}{Breakpoints and damped restart}
A pulse edge inside a step is invisible: the companion history belongs to
the pre-edge circuit. For $h/\tau = 0.1$:
\[
v_1 = \frac{h/2\tau}{1 + h/2\tau} \approx 0.048
\quad\text{instead of}\quad
1 - e^{-h/\tau} \approx 0.095
\]
a factor-two error no smaller $h$ fixes.
\vspace{2mm}
\begin{itemize}
\item every point of every piecewise linear waveform is a breakpoint;
      steps land on them exactly
\item after a breakpoint: small step, cleared predictor history
\item the two predictor-less bootstrap steps run backward Euler: first
      order but strongly damped, absorbing the trapezoid ringing at the
      discontinuity
\end{itemize}
\end{frame}

\begin{frame}[fragile]{Small-signal AC}
\[
(G + j\omega C)\,\hat{x} = \hat{b}
\]
At a converged operating point the limiters pass voltages through
unchanged, so $G$ \emph{is exactly the DC Jacobian the last Newton
iteration assembled}. The AC path reuses \texttt{assemble}, maps to
complex, adds $j\omega C$ and $-j\omega L$.
\begin{rust}
trait Num: Copy + Add<Output = Self> + Sub<Output = Self>
         + Mul<Output = Self> + Div<Output = Self> {
    fn zero() -> Self;
    fn one() -> Self;
    fn mag(self) -> f64;   // pivot magnitude: abs for f64, modulus for Cx
}
\end{rust}
One generic sparse LU serves Newton and AC; no device model contains any
AC-specific code.
\end{frame}

\begin{frame}[fragile]{It oscillates}
\begin{columns}[T]
\column{0.52\textwidth}
\begin{rust}
cmos ring oscillator
.model n nmos kp=1e-3 vt0=1 cgs=5p cgd=2p
.model p pmos kp=1e-3 vt0=1 cgs=5p cgd=2p
vdd vdd 0 dc 5
m1 b a 0 0 n
m2 b a vdd vdd p
m3 c b 0 0 n
m4 c b vdd vdd p
m5 a c 0 0 n
m6 a c vdd vdd p
i1 0 a pulse(0 1m 0 1n 1n 20n)
.tran 1n 500n
\end{rust}
\column{0.44\textwidth}
\begin{itemize}
\item three inverters, gate capacitances set the stage delay
\item current kick leaves the metastable operating point
\item period 78\,ns, stable within 1.5\,\%
\item a test measures the period and bounds the jitter
\end{itemize}
\vspace{2mm}
Without the capacitances the transistors switch instantaneously and the
ring has no defined frequency.
\end{columns}
\end{frame}

\begin{frame}{Validation}
\locbar{tests/cli.rs}{330 lines, 0 of 992: tests are free}{0}{0}
\vspace{1mm}
References independent of the implementation:
\begin{itemize}
\item analytic: RC/RL step vs. $1 - e^{-t/\tau}$, AC corner at $1/\sqrt{2}$
      and $-45^\circ$, LC amplitude $i_0\sqrt{L/C}$, energy conservation
\item bias points: MOSFET, JFET, BJT ($I_C = \beta_F I_B$), exact npn/pnp
      symmetry
\item structural: a randomized 30-stage ladder verified against a Thevenin
      reduction computed independently in the test
\item negative: seven malformed decks must produce named errors, never a
      panic; the floating node error must name the node
\item meta: the line counter itself
\end{itemize}
\end{frame}

\begin{frame}{The cut list}
\begin{center}
\begin{tabular}{llr}
\toprule
feature & why it lost & est. lines \\
\midrule
\texttt{.subckt} hierarchy & lost to the transistor models & $\sim$100 \\
Markowitz ordering & netlist order suffices at this scale & $\sim$80 \\
Gummel-Poon BJT & Ebers-Moll covers the teaching core & $\sim$80 \\
noise analysis & needs the adjoint system & $\sim$80 \\
\texttt{.param} expressions & a calculator is its own project & $\sim$60 \\
temperature sweeps & $V_T$ is a constant, by choice & $\sim$25 \\
\bottomrule
\end{tabular}
\end{center}
\vspace{2mm}
The most consequential remaining cut is \texttt{.subckt}: real decks are
hierarchical. It would be the first addition if the budget grew.
\end{frame}

\begin{frame}{Closing}
\begin{center}
{\large \textbf{992 of 1000 lines}}
\end{center}
\vspace{2mm}
MNA, sparse LU with partial pivoting, two contribution primitives
carrying all device stamps, Newton with pnjlim and data-driven homotopy
fallbacks, trapezoidal integration with LTE control, breakpoints and
damped restarts, charge and flux based reactive state with graded
junction capacitances, small-signal AC, eleven device letters, four
analyses, 23 tests.
\vspace{4mm}
\begin{center}
\texttt{github.com/milanofthe/nanospice}\\[1mm]
{\footnotesize MIT licensed; the design report with all derivations is in
the repository.}
\end{center}
\end{frame}

\end{document}
